\section*{Introdução}

A intensificação das mudanças climáticas tem provocado o aumento da frequência e da severidade de desastres ambientais, como incêndios florestais. Em áreas de proteção ambiental e segurança civil, o mero registro isolado de um foco de calor em um mapa é insuficiente para determinar sua verdadeira ameaça. Para que a informação seja estratégica e embase ações governamentais, ela deve ser correlacionada com as vulnerabilidades do território afetado. Diante dessa necessidade, a questão-problema central que este estudo investiga é: como orquestrar e cruzar interativamente grandes massas dinâmicas de telemetria de incêndios com infraestruturas críticas, superando os limites de latência dos sistemas tradicionais para prover consciência situacional efetiva?

O objetivo deste trabalho é apresentar a CCO-DataOps, uma plataforma baseada em Engenharia de Dados orientada a eventos para fornecer consciência situacional em sub-segundos. A plataforma foca na integração contínua de anomalias térmicas de satélite com infraestruturas críticas e aspectos socioambientais, utilizando o estado de Minas Gerais como base de validação técnica. Ao mapear o impacto primário e as distâncias até hospitais, malhas viárias e terras indígenas, o projeto demonstra a viabilidade de aplicar alto processamento analítico para subsidiar respostas de emergência.
